% Options for packages loaded elsewhere
\PassOptionsToPackage{unicode}{hyperref}
\PassOptionsToPackage{hyphens}{url}
\documentclass[
]{article}
\usepackage{xcolor}
\usepackage[margin=1in]{geometry}
\usepackage{amsmath,amssymb}
\setcounter{secnumdepth}{-\maxdimen} % remove section numbering
\usepackage{iftex}
\ifPDFTeX
  \usepackage[T1]{fontenc}
  \usepackage[utf8]{inputenc}
  \usepackage{textcomp} % provide euro and other symbols
\else % if luatex or xetex
  \usepackage{unicode-math} % this also loads fontspec
  \defaultfontfeatures{Scale=MatchLowercase}
  \defaultfontfeatures[\rmfamily]{Ligatures=TeX,Scale=1}
\fi
\usepackage{lmodern}
\ifPDFTeX\else
  % xetex/luatex font selection
\fi
% Use upquote if available, for straight quotes in verbatim environments
\IfFileExists{upquote.sty}{\usepackage{upquote}}{}
\IfFileExists{microtype.sty}{% use microtype if available
  \usepackage[]{microtype}
  \UseMicrotypeSet[protrusion]{basicmath} % disable protrusion for tt fonts
}{}
\makeatletter
\@ifundefined{KOMAClassName}{% if non-KOMA class
  \IfFileExists{parskip.sty}{%
    \usepackage{parskip}
  }{% else
    \setlength{\parindent}{0pt}
    \setlength{\parskip}{6pt plus 2pt minus 1pt}}
}{% if KOMA class
  \KOMAoptions{parskip=half}}
\makeatother
\usepackage{color}
\usepackage{fancyvrb}
\newcommand{\VerbBar}{|}
\newcommand{\VERB}{\Verb[commandchars=\\\{\}]}
\DefineVerbatimEnvironment{Highlighting}{Verbatim}{commandchars=\\\{\}}
% Add ',fontsize=\small' for more characters per line
\usepackage{framed}
\definecolor{shadecolor}{RGB}{248,248,248}
\newenvironment{Shaded}{\begin{snugshade}}{\end{snugshade}}
\newcommand{\AlertTok}[1]{\textcolor[rgb]{0.94,0.16,0.16}{#1}}
\newcommand{\AnnotationTok}[1]{\textcolor[rgb]{0.56,0.35,0.01}{\textbf{\textit{#1}}}}
\newcommand{\AttributeTok}[1]{\textcolor[rgb]{0.13,0.29,0.53}{#1}}
\newcommand{\BaseNTok}[1]{\textcolor[rgb]{0.00,0.00,0.81}{#1}}
\newcommand{\BuiltInTok}[1]{#1}
\newcommand{\CharTok}[1]{\textcolor[rgb]{0.31,0.60,0.02}{#1}}
\newcommand{\CommentTok}[1]{\textcolor[rgb]{0.56,0.35,0.01}{\textit{#1}}}
\newcommand{\CommentVarTok}[1]{\textcolor[rgb]{0.56,0.35,0.01}{\textbf{\textit{#1}}}}
\newcommand{\ConstantTok}[1]{\textcolor[rgb]{0.56,0.35,0.01}{#1}}
\newcommand{\ControlFlowTok}[1]{\textcolor[rgb]{0.13,0.29,0.53}{\textbf{#1}}}
\newcommand{\DataTypeTok}[1]{\textcolor[rgb]{0.13,0.29,0.53}{#1}}
\newcommand{\DecValTok}[1]{\textcolor[rgb]{0.00,0.00,0.81}{#1}}
\newcommand{\DocumentationTok}[1]{\textcolor[rgb]{0.56,0.35,0.01}{\textbf{\textit{#1}}}}
\newcommand{\ErrorTok}[1]{\textcolor[rgb]{0.64,0.00,0.00}{\textbf{#1}}}
\newcommand{\ExtensionTok}[1]{#1}
\newcommand{\FloatTok}[1]{\textcolor[rgb]{0.00,0.00,0.81}{#1}}
\newcommand{\FunctionTok}[1]{\textcolor[rgb]{0.13,0.29,0.53}{\textbf{#1}}}
\newcommand{\ImportTok}[1]{#1}
\newcommand{\InformationTok}[1]{\textcolor[rgb]{0.56,0.35,0.01}{\textbf{\textit{#1}}}}
\newcommand{\KeywordTok}[1]{\textcolor[rgb]{0.13,0.29,0.53}{\textbf{#1}}}
\newcommand{\NormalTok}[1]{#1}
\newcommand{\OperatorTok}[1]{\textcolor[rgb]{0.81,0.36,0.00}{\textbf{#1}}}
\newcommand{\OtherTok}[1]{\textcolor[rgb]{0.56,0.35,0.01}{#1}}
\newcommand{\PreprocessorTok}[1]{\textcolor[rgb]{0.56,0.35,0.01}{\textit{#1}}}
\newcommand{\RegionMarkerTok}[1]{#1}
\newcommand{\SpecialCharTok}[1]{\textcolor[rgb]{0.81,0.36,0.00}{\textbf{#1}}}
\newcommand{\SpecialStringTok}[1]{\textcolor[rgb]{0.31,0.60,0.02}{#1}}
\newcommand{\StringTok}[1]{\textcolor[rgb]{0.31,0.60,0.02}{#1}}
\newcommand{\VariableTok}[1]{\textcolor[rgb]{0.00,0.00,0.00}{#1}}
\newcommand{\VerbatimStringTok}[1]{\textcolor[rgb]{0.31,0.60,0.02}{#1}}
\newcommand{\WarningTok}[1]{\textcolor[rgb]{0.56,0.35,0.01}{\textbf{\textit{#1}}}}
\usepackage{graphicx}
\makeatletter
\newsavebox\pandoc@box
\newcommand*\pandocbounded[1]{% scales image to fit in text height/width
  \sbox\pandoc@box{#1}%
  \Gscale@div\@tempa{\textheight}{\dimexpr\ht\pandoc@box+\dp\pandoc@box\relax}%
  \Gscale@div\@tempb{\linewidth}{\wd\pandoc@box}%
  \ifdim\@tempb\p@<\@tempa\p@\let\@tempa\@tempb\fi% select the smaller of both
  \ifdim\@tempa\p@<\p@\scalebox{\@tempa}{\usebox\pandoc@box}%
  \else\usebox{\pandoc@box}%
  \fi%
}
% Set default figure placement to htbp
\def\fps@figure{htbp}
\makeatother
\setlength{\emergencystretch}{3em} % prevent overfull lines
\providecommand{\tightlist}{%
  \setlength{\itemsep}{0pt}\setlength{\parskip}{0pt}}
\usepackage{bookmark}
\IfFileExists{xurl.sty}{\usepackage{xurl}}{} % add URL line breaks if available
\urlstyle{same}
\hypersetup{
  hidelinks,
  pdfcreator={LaTeX via pandoc}}

\author{}
\date{\vspace{-2.5em}}

\begin{document}

\section{Extraction de connaissances à partir de données structurées et
non
structurées}\label{extraction-de-connaissances-uxe0-partir-de-donnuxe9es-structuruxe9es-et-non-structuruxe9es}

\subsection{Séance finale : TP noté
2}\label{suxe9ance-finale-tp-notuxe9-2}

Déposer le fichier \texttt{.ipynb} produit (attention à indiquer votre
nom de famille dans le nom de fichier) dans ce dépôt :

\url{https://cloud.parisdescartes.fr/index.php/s/6g9zyxZjkPwDB6z}

Délai : A rendre \textbf{avant la fin de la séance} le jour même

\subsection{Données}\label{donnuxe9es}

Nous reprenons le jeu de données du premier TP noté, contenant des
informations sur des menus McDonald. Vous trouverez le détail des
variables sur
\href{https://www.kaggle.com/mcdonalds/nutrition-facts}{cette page}.
Vous devez l'importer comme ci-dessous.

\begin{Shaded}
\begin{Highlighting}[]
\ImportTok{import}\NormalTok{ pandas}

\NormalTok{menus }\OperatorTok{=}\NormalTok{ pandas.read\_csv(}\StringTok{"https://fxjollois.github.io/cours{-}2025{-}2026/ufr{-}{-}m1{-}dci{-}{-}ecd/menus\_mcdonald.csv"}\NormalTok{)}
\NormalTok{menus.head()}
\end{Highlighting}
\end{Shaded}

Unnamed: 0

Category

Item

Serving Size

Calories

Calories from Fat

Total Fat

Total Fat (\% Daily Value)

Saturated Fat

Saturated Fat (\% Daily Value)

\ldots{}

Carbohydrates

Carbohydrates (\% Daily Value)

Dietary Fiber

Dietary Fiber (\% Daily Value)

Sugars

Protein

Vitamin A (\% Daily Value)

Vitamin C (\% Daily Value)

Calcium (\% Daily Value)

Iron (\% Daily Value)

0

0

Breakfast

Egg McMuffin

4.8 oz (136 g)

300

120

13.0

20

5.0

25

\ldots{}

31

10

4

17

3

17

10

0

25

15

1

1

Breakfast

Egg White Delight

4.8 oz (135 g)

250

70

8.0

12

3.0

15

\ldots{}

30

10

4

17

3

18

6

0

25

8

2

2

Breakfast

Sausage McMuffin

3.9 oz (111 g)

370

200

23.0

35

8.0

42

\ldots{}

29

10

4

17

2

14

8

0

25

10

3

3

Breakfast

Sausage McMuffin with Egg

5.7 oz (161 g)

450

250

28.0

43

10.0

52

\ldots{}

30

10

4

17

2

21

15

0

30

15

4

4

Breakfast

Sausage McMuffin with Egg Whites

5.7 oz (161 g)

400

210

23.0

35

8.0

42

\ldots{}

30

10

4

17

2

21

6

0

25

10

5 rows × 25 columns

\subsection{A faire}\label{a-faire}

Sous la forme d'un notebook le plus propre et le mieux organisé
possible, vous devez :

\begin{itemize}
\tightlist
\item
  Décrire les données (distribution des variables, données aberrantes)
\item
  Réperer les variables sur lesquelles se concentrer (il y en a
  plusieurs à supprimer)
\item
  Réaliser une ACP sur les données

  \begin{itemize}
  \tightlist
  \item
    en justifiant s'il faut ou non standardiser les variables au
    préalable
  \end{itemize}
\item
  Représenter les menus sur le plan factoriel

  \begin{itemize}
  \tightlist
  \item
    commenter celui-ci
  \end{itemize}
\item
  Décrire les axes en se basant sur le cercle des corrélations des
  variables
\item
  Déterminer un nombre de classes adapté à ce jeu de données
\item
  Trouver la partition des menus, avec le nombre de classes optimal
  obtenu ci-dessus
\item
  Décrire les classes et les comparer entre elles
\item
  Effectuer toutes les représentations qui vous semblent utiles pour
  décrire vos résultats et accompagner votre conclusion
\end{itemize}

\end{document}
